% GENERATED FILE. Edit canonical lessons, then run npm run generate:books.
% canonical-source-hash: fnv1a64:1df4c1f99a08964c
% canonical-lessons: UR-C516-khattati, UR-C516-naqqashi, UR-C516-murti, UR-C516-kanda-karna, UR-C516-but

\chapter{Carving and Calligraphy}
\label{ch:ur-carving-and-calligraphy}
\hlchaptermodality{516}

\begin{chapteropening}
\textbf{By the end of this chapter:} \emph{I can talk about beautiful writing and carving.}

Complete the last lesson of chapter 516: I can talk about beautiful writing and carving.
\end{chapteropening}

\section[\texorpdfstring{\ur{خطاطی} (khattātī)}{خطاطی (khattātī)}]{\ur{خطاطی} (khattātī) — calligraphy}

\label{lesson:UR-C516-khattati}



\begin{quote}
Before the new one: say the Urdu for a masterpiece, then the Urdu for to show through.
\end{quote}

\subsection*{You'll want to know: \ur{خطاطی}}
\textbf{\ur{خطاطی}} — \emph{khattātī} — \textquotedblleft{}calligraphy\textquotedblright{}.

\begin{quote}
\ur{میرے} \ur{استاد} \ur{خطاطی} \ur{سکھاتے} \ur{ہیں۔}
\end{quote}

\begin{quote}
\emph{mere ustād khattātī sikhāte haiṅ.}
\end{quote}

\begin{quote}
My teacher teaches calligraphy.
\end{quote}

Feminine.

\begin{cousinweb}
\textbf{\ur{خطاطی}} (\emph{khattātī}) is built on \emph{khattāt}, \textquotedblleft{}a calligrapher\textquotedblright{}, usually traced, through Persian, to Arabic \emph{khaṭṭāṭ}, with the noun ending \emph{-ī}.
\end{cousinweb}

\begin{grammarlens}[title={What you've built}]
One more word for describing what you lived and felt.
\end{grammarlens}

\subsection*{Your turn}
\begin{itemize}
\raggedright
  \item \emph{Say it:} \emph{khattātī}
  \item \emph{Say it:} \emph{khattātī}, once more
  \item \emph{Say it:} \emph{mere ustād khattātī sikhāte haiṅ.} — My teacher teaches calligraphy.
  \item \emph{Recall:} say the Urdu for a masterpiece, then the Urdu for to show through, then say \emph{khattātī} again
\end{itemize}

\begin{tcolorbox}[breakable,colback=teal!4,colframe=teal!35!black,title={Before you move on}]
What does \ur{خطاطی} mean? (\textquotedblleft{}Calligraphy\textquotedblright{}.) Say it once more.
\end{tcolorbox}
\section[\texorpdfstring{\ur{نقاشی} (naqqāshī)}{نقاشی (naqqāshī)}]{\ur{نقاشی} (naqqāshī) — carved or painted decoration}

\label{lesson:UR-C516-naqqashi}



\begin{quote}
Before the new one: say the Urdu for to show through, then the Urdu for calligraphy.
\end{quote}

\subsection*{You'll want to know: \ur{نقاشی}}
\textbf{\ur{نقاشی}} — \emph{naqqāshī} — \textquotedblleft{}carved or painted decoration\textquotedblright{}.

\begin{quote}
\ur{اس} \ur{پرانے} \ur{دروازے} \ur{پر} \ur{نقاشی} \ur{ہے۔}
\end{quote}

\begin{quote}
\emph{is purāne darvāze par naqqāshī hai.}
\end{quote}

\begin{quote}
This old door has carving all over it.
\end{quote}

Feminine.

\begin{cousinweb}
\textbf{\ur{نقاشی}} (\emph{naqqāshī}) is built on \emph{naqqāsh}, \textquotedblleft{}a carver, a painter\textquotedblright{}, usually traced, through Persian, to Arabic \emph{naqqāsh}, with the noun ending \emph{-ī}.
\end{cousinweb}

\begin{grammarlens}[title={What you've built}]
One more word for describing what you lived and felt.
\end{grammarlens}

\subsection*{Your turn}
\begin{itemize}
\raggedright
  \item \emph{Say it:} \emph{naqqāshī}
  \item \emph{Say it:} \emph{naqqāshī}, once more
  \item \emph{Say it:} \emph{is purāne darvāze par naqqāshī hai.} — This old door has carving all over it.
  \item \emph{Recall:} say the Urdu for to show through, then the Urdu for calligraphy, then say \emph{naqqāshī} again
\end{itemize}

\begin{tcolorbox}[breakable,colback=teal!4,colframe=teal!35!black,title={Before you move on}]
What does \ur{نقاشی} mean? (\textquotedblleft{}Carved or painted decoration\textquotedblright{}.) Say it once more.
\end{tcolorbox}
\section[\texorpdfstring{\ur{مورتی} (mūrtī)}{مورتی (mūrtī)}]{\ur{مورتی} (mūrtī) — a figurine; a small statue}

\label{lesson:UR-C516-murti}



\begin{quote}
Before the new one: say the Urdu for calligraphy, then the Urdu for carved or painted decoration.
\end{quote}

\subsection*{You'll want to know: \ur{مورتی}}
\textbf{\ur{مورتی}} — \emph{mūrtī} — \textquotedblleft{}a figurine; a small statue\textquotedblright{}.

\begin{quote}
\ur{کمہار} \ur{نے} \ur{مٹی} \ur{کی} \ur{ایک} \ur{مورتی} \ur{بنائی۔}
\end{quote}

\begin{quote}
\emph{kumhār ne miṭṭī kī ek mūrtī banāī.}
\end{quote}

\begin{quote}
The potter made a little clay figure.
\end{quote}

Feminine.

\begin{cousinweb}
\textbf{\ur{مورتی}} (\emph{mūrtī}) is inherited from Sanskrit \emph{mūrti}, \textquotedblleft{}a solid form; an image\textquotedblright{} (Monier-Williams).
\end{cousinweb}

\begin{grammarlens}[title={What you've built}]
One more word for describing what you lived and felt.
\end{grammarlens}

\subsection*{Your turn}
\begin{itemize}
\raggedright
  \item \emph{Say it:} \emph{mūrtī}
  \item \emph{Say it:} \emph{mūrtī}, once more
  \item \emph{Say it:} \emph{kumhār ne miṭṭī kī ek mūrtī banāī.} — The potter made a little clay figure.
  \item \emph{Recall:} say the Urdu for calligraphy, then the Urdu for carved or painted decoration, then say \emph{mūrtī} again
\end{itemize}

\begin{tcolorbox}[breakable,colback=teal!4,colframe=teal!35!black,title={Before you move on}]
What does \ur{مورتی} mean? (\textquotedblleft{}A figurine; a small statue\textquotedblright{}.) Say it once more.
\end{tcolorbox}
\section[\texorpdfstring{\ur{کندہ} \ur{کرنا} (kandā karnā)}{کندہ کرنا (kandā karnā)}]{\ur{کندہ} \ur{کرنا} (kandā karnā) — to engrave}

\label{lesson:UR-C516-kanda-karna}



\begin{quote}
Before the new one: say the Urdu for carved or painted decoration, then the Urdu for a figurine.
\end{quote}

\subsection*{You'll want to know: \ur{کندہ} \ur{کرنا}}
\textbf{\ur{کندہ} \ur{کرنا}} — \emph{kandā karnā} — \textquotedblleft{}to engrave\textquotedblright{}.

\begin{quote}
\ur{سنار} \ur{نے} \ur{انگوٹھی} \ur{پر} \ur{نام} \ur{کندہ} \ur{کیا۔}
\end{quote}

\begin{quote}
\emph{sunār ne angūṭhī par nām kandā kiyā.}
\end{quote}

\begin{quote}
The goldsmith engraved a name on the ring.
\end{quote}

Past, with \ur{نے} (he or she …): \ur{اس} \ur{نے} \ur{پتھر} \ur{پر} \ur{ایک} \ur{لفظ} \ur{کندہ} \ur{کیا}, \emph{us ne patthar par ek lafz kandā kiyā}.

\begin{cousinweb}
\textbf{\ur{کندہ} \ur{کرنا}} (\emph{kandā karnā}) is built from \emph{kandā}, usually traced to Persian \emph{kanda}, \textquotedblleft{}dug, carved\textquotedblright{}, from \emph{kandan}, \textquotedblleft{}to dig\textquotedblright{}, and \emph{karnā}, \textquotedblleft{}to do\textquotedblright{}.
\end{cousinweb}

\begin{grammarlens}[title={What you've built}]
One more word for describing what you lived and felt.
\end{grammarlens}

\subsection*{Your turn}
\begin{itemize}
\raggedright
  \item \emph{Say it:} \emph{kandā karnā}
  \item \emph{Say it:} \emph{kandā karnā}, once more
  \item \emph{Say it:} \emph{sunār ne angūṭhī par nām kandā kiyā.} — The goldsmith engraved a name on the ring.
  \item \emph{Recall:} say the Urdu for carved or painted decoration, then the Urdu for a figurine, then say \emph{kandā karnā} again
\end{itemize}

\begin{tcolorbox}[breakable,colback=teal!4,colframe=teal!35!black,title={Before you move on}]
What does \ur{کندہ} \ur{کرنا} mean? (\textquotedblleft{}To engrave\textquotedblright{}.) Say it once more.
\end{tcolorbox}
\section[\texorpdfstring{\ur{بت} (but)}{بت (but)}]{\ur{بت} (but) — an idol; a carved figure}

\label{lesson:UR-C516-but}



\begin{quote}
Before the new one: say the Urdu for a figurine, then the Urdu for to engrave.
\end{quote}

\subsection*{You'll want to know: \ur{بت}}
\textbf{\ur{بت}} — \emph{but} — \textquotedblleft{}an idol; a carved figure\textquotedblright{}.

\begin{quote}
\ur{اس} \ur{نے} \ur{پتھر} \ur{سے} \ur{ایک} \ur{بت} \ur{بنایا۔}
\end{quote}

\begin{quote}
\emph{us ne patthar se ek but banāyā.}
\end{quote}

\begin{quote}
He carved a figure out of stone.
\end{quote}

Masculine.

\begin{cousinweb}
\textbf{\ur{بت}} (\emph{but}) is usually traced to Persian \emph{but}, \textquotedblleft{}an idol, a carved figure\textquotedblright{}.
\end{cousinweb}

\begin{grammarlens}[title={What you've built}]
One more word for describing what you lived and felt.
\end{grammarlens}

\subsection*{Your turn}
\begin{itemize}
\raggedright
  \item \emph{Say it:} \emph{but}
  \item \emph{Say it:} \emph{but}, once more
  \item \emph{Say it:} \emph{us ne patthar se ek but banāyā.} — He carved a figure out of stone.
  \item \emph{Recall:} say the Urdu for a figurine, then the Urdu for to engrave, then say \emph{but} again
\end{itemize}

\begin{tcolorbox}[breakable,colback=teal!4,colframe=teal!35!black,title={Before you move on}]
What does \ur{بت} mean? (\textquotedblleft{}An idol; a carved figure\textquotedblright{}.) Say it once more.
\end{tcolorbox}
